% ============================================================
%  Docker Dependency Specification – Landfill e-Forms Application
%  City of Los Angeles, Department of Sanitation
%  Cal State LA Senior Design 2017–2018
% ============================================================
\documentclass[12pt,letterpaper]{article}

% ── Packages ────────────────────────────────────────────────
\usepackage[margin=1in]{geometry}
\usepackage{titlesec}
\usepackage{hyperref}
\usepackage{enumitem}
\usepackage{booktabs}
\usepackage{longtable}
\usepackage{array}
\usepackage{xcolor}
\usepackage{fancyhdr}
\usepackage{parskip}
\usepackage{listings}
\usepackage{microtype}
\usepackage{colortbl}
\usepackage{cellspace}

% ── Colours ─────────────────────────────────────────────────
\definecolor{headerblue}{HTML}{1F5C99}
\definecolor{subblue}{HTML}{2E75B6}
\definecolor{rowshade}{HTML}{EFF4FB}
\definecolor{codebg}{HTML}{F4F4F4}
\definecolor{muted}{HTML}{555555}
\definecolor{tablerule}{HTML}{CCCCCC}

% ── Hyperlinks ───────────────────────────────────────────────
\hypersetup{
  colorlinks=true,
  linkcolor=headerblue,
  urlcolor=headerblue
}

% ── Header / Footer ─────────────────────────────────────────
\pagestyle{fancy}
\fancyhf{}
\lhead{\color{muted}\small Landfill e-Forms Application\ |\ Docker Dependency Specification}
\cfoot{\color{muted}\small\thepage}
\renewcommand{\headrulewidth}{0.4pt}
\renewcommand{\footrulewidth}{0pt}

% ── Section formatting ───────────────────────────────────────
\titleformat{\section}
  {\large\bfseries\color{headerblue}}
  {}
  {0em}
  {}
  [\color{headerblue}\titlerule]

\titleformat{\subsection}
  {\normalsize\bfseries\color{subblue}}
  {}
  {0em}
  {}

% ── Table helpers ────────────────────────────────────────────
\setlength{\tabcolsep}{8pt}
\renewcommand{\arraystretch}{1.35}

\newcommand{\hdrrow}[1]{%
  \rowcolor{headerblue}%
  \color{white}\textbf{#1}%
}

% ── Code listing style ───────────────────────────────────────
\lstset{
  basicstyle=\ttfamily\small,
  backgroundcolor=\color{codebg},
  frame=single,
  framerule=0.4pt,
  rulecolor=\color{tablerule},
  breaklines=true,
  captionpos=b,
  xleftmargin=0.5cm,
  xrightmargin=0.5cm
}

% ============================================================
\begin{document}

% ── Title Block ──────────────────────────────────────────────
\begin{center}
  {\Huge\bfseries\color{headerblue} Landfill e-Forms Application}\\[0.5em]
  {\Large\color{subblue} Docker Dependency Specification}\\[0.8em]
  {\small\color{muted}
    City of Los Angeles, Department of Sanitation
    \enspace$\cdot$\enspace
    Cal State LA Senior Design 2017--2018}
\end{center}

\vspace{1em}
\hrule height 0.8pt
\vspace{1.5em}



% ============================================================
\section{Docker Images}

Four Docker images must be pulled before building or running any container.

\vspace{0.5em}
\begin{longtable}{>{\ttfamily}p{4.2cm} p{5.8cm} p{3.2cm}}
  \rowcolor{headerblue}
  \color{white}\textbf{Image} &
  \color{white}\textbf{Role} &
  \color{white}\textbf{Source} \\
  \endfirsthead
  \rowcolor{headerblue}
  \color{white}\textbf{Image} &
  \color{white}\textbf{Role} &
  \color{white}\textbf{Source} \\
  \endhead
  openjdk:8-jdk-alpine  & Spring Boot back-end JAR runtime          & Docker Hub \\
  \rowcolor{rowshade}
  mysql:5.7             & Relational database server                 & Docker Hub \\
  node:6-alpine         & Angular 4 front-end build stage            & Docker Hub \\
  \rowcolor{rowshade}
  nginx:alpine          & Serves compiled Angular static bundle      & Docker Hub \\
\end{longtable}

% ============================================================
\section{Back-End Container \textnormal{\small(\texttt{openjdk:8-jdk-alpine})}}

All Java dependencies are resolved at build time by Maven from the project
\texttt{pom.xml}. The resulting fat-JAR bundles everything listed below —
no manual installation is required inside the running container.

\vspace{0.5em}
\begin{longtable}{p{5.5cm} p{1.8cm} p{6cm}}
  \rowcolor{headerblue}
  \color{white}\textbf{Dependency (artifactId)} &
  \color{white}\textbf{Version} &
  \color{white}\textbf{Purpose} \\
  \endfirsthead
  \rowcolor{headerblue}
  \color{white}\textbf{Dependency (artifactId)} &
  \color{white}\textbf{Version} &
  \color{white}\textbf{Purpose} \\
  \endhead
  spring-boot-starter-web       & 2.x    & Spring MVC REST controllers + embedded Tomcat \\
  \rowcolor{rowshade}
  spring-boot-starter-data-jpa  & 2.x    & Spring Data JPA repository layer \\
  spring-boot-starter-security  & 2.x    & Authentication \& role-based authorisation \\
  \rowcolor{rowshade}
  hibernate-core                & 5.x    & JPA implementation / ORM / SQL generation \\
  mysql-connector-java          & 5.1.x  & JDBC driver for MySQL 5.7 \\
  \rowcolor{rowshade}
  jjwt                          & 0.9.x  & JSON Web Token creation \& validation \\
  jasperreports                 & 6.x    & PDF \& CSV report generation \\
  \rowcolor{rowshade}
  spring-boot-maven-plugin      & 2.x    & Packages application as executable fat-JAR \\
  maven-compiler-plugin         & 3.x    & Compiles Java 8 source \\
  \rowcolor{rowshade}
  lombok                        & 1.18.x & Boilerplate reduction (getters/setters/builders) \\
\end{longtable}

\vspace{0.5em}
Build command executed inside the container build stage:

\begin{lstlisting}[language=bash]
mvn clean package -DskipTests
\end{lstlisting}

% ============================================================
\section{Front-End Build Container \textnormal{\small(\texttt{node:6-alpine})}}

\texttt{npm install} resolves all dependencies listed in \texttt{package.json}.
The compiled static bundle is then copied into the Nginx container; the Node
container is discarded after the build stage.

\vspace{0.5em}
\begin{longtable}{p{5.5cm} p{1.8cm} p{6cm}}
  \rowcolor{headerblue}
  \color{white}\textbf{Package} &
  \color{white}\textbf{Version} &
  \color{white}\textbf{Purpose} \\
  \endfirsthead
  \rowcolor{headerblue}
  \color{white}\textbf{Package} &
  \color{white}\textbf{Version} &
  \color{white}\textbf{Purpose} \\
  \endhead
  @angular/core              & 4.x   & Core Angular framework \\
  \rowcolor{rowshade}
  @angular/common            & 4.x   & Common directives \& pipes \\
  @angular/compiler          & 4.x   & Template compiler \\
  \rowcolor{rowshade}
  @angular/forms             & 4.x   & Reactive \& template-driven forms \\
  @angular/http              & 4.x   & HTTP client module \\
  \rowcolor{rowshade}
  @angular/router            & 4.x   & Client-side routing \\
  @angular/platform-browser  & 4.x   & Browser rendering platform \\
  \rowcolor{rowshade}
  rxjs                       & 5.x   & Reactive extensions / observable HTTP calls \\
  bootstrap                  & 3.x   & Responsive CSS grid \& UI components \\
  \rowcolor{rowshade}
  zone.js                    & 0.8.x & Angular change-detection runtime \\
  typescript                 & 2.x   & Typed superset of JavaScript \\
  \rowcolor{rowshade}
  @angular/cli               & 1.x   & Build toolchain, dev server, \texttt{ng build -{}-prod} \\
  @angular/compiler-cli      & 4.x   & Ahead-of-time (AOT) compiler \\
  \rowcolor{rowshade}
  webpack                    & 2.x   & Module bundler (used internally by Angular CLI) \\
\end{longtable}

\vspace{0.5em}
Build command executed inside the Node build stage:

\begin{lstlisting}[language=bash]
npm install && ng build --prod --aot
\end{lstlisting}

% ============================================================
\section{Web-Server Container \textnormal{\small(\texttt{nginx:alpine})}}

No additional packages are installed inside this container. The following
configuration is applied:

\begin{itemize}[noitemsep]
  \item \texttt{nginx:alpine} base image
  \item Compiled Angular bundle copied into \texttt{/usr/share/nginx/html}
  \item Custom \texttt{nginx.conf} to proxy \texttt{/api/*} requests to the
        Spring Boot back-end container
  \item \texttt{try\_files \$uri \$uri/ /index.html} for Angular client-side routing
\end{itemize}

% ============================================================
\section{Database Container \textnormal{\small(\texttt{mysql:5.7})}}

No additional packages are installed inside this container. Initialisation
is handled via Docker volume mounts:

\begin{itemize}[noitemsep]
  \item \texttt{mysql:5.7} base image — no additional packages required
  \item \texttt{schema.sql} mounted as a Docker volume for database
        initialisation on first run
  \item Environment variables required: \texttt{MYSQL\_ROOT\_PASSWORD},
        \texttt{MYSQL\_DATABASE}, \texttt{MYSQL\_USER}, \texttt{MYSQL\_PASSWORD}
\end{itemize}

% ============================================================
\section{Android Mobile Application \textnormal{\small(not containerised)}}

The Android app runs on the inspector's device, not inside Docker. It is
built locally via Gradle. The following dependencies are declared in
\texttt{build.gradle} and downloaded from Maven Central / Google Maven at
build time:

\vspace{0.5em}
\begin{longtable}{p{5.5cm} p{1.8cm} p{6cm}}
  \rowcolor{headerblue}
  \color{white}\textbf{Dependency} &
  \color{white}\textbf{Version} &
  \color{white}\textbf{Purpose} \\
  \endfirsthead
  \rowcolor{headerblue}
  \color{white}\textbf{Dependency} &
  \color{white}\textbf{Version} &
  \color{white}\textbf{Purpose} \\
  \endhead
  appcompat-v7                    & 27.x & Material UI components \& back-compat \\
  \rowcolor{rowshade}
  support:design                  & 27.x & Navigation drawer, FAB, Snackbar \\
  retrofit2:retrofit              & 2.x  & Type-safe REST HTTP client \\
  \rowcolor{rowshade}
  retrofit2:converter-gson        & 2.x  & JSON $\to$ Java object conversion \\
  gson                            & 2.x  & JSON serialisation / deserialisation \\
  \rowcolor{rowshade}
  room:runtime                    & 1.x  & SQLite ORM (offline local storage) \\
  room:compiler                   & 1.x  & Room annotation processor \\
  \rowcolor{rowshade}
  okhttp3:okhttp                  & 3.x  & Underlying HTTP engine for Retrofit \\
  okhttp3:logging-interceptor     & 3.x  & HTTP request/response logging \\
  \rowcolor{rowshade}
  junit                           & 4.x  & Unit testing framework \\
\end{longtable}

% ============================================================
\section{docker-compose.yml — Service Summary}

The four Docker services and their inter-dependencies:

\vspace{0.5em}
\begin{longtable}{p{3cm} p{4.5cm} p{2.8cm} p{2.8cm}}
  \rowcolor{headerblue}
  \color{white}\textbf{Service} &
  \color{white}\textbf{Image} &
  \color{white}\textbf{Depends On} &
  \color{white}\textbf{Port} \\
  \endfirsthead
  \rowcolor{headerblue}
  \color{white}\textbf{Service} &
  \color{white}\textbf{Image} &
  \color{white}\textbf{Depends On} &
  \color{white}\textbf{Port} \\
  \endhead
  \texttt{db}               & \texttt{mysql:5.7}                    & —              & 3306 (internal) \\
  \rowcolor{rowshade}
  \texttt{backend}          & \texttt{openjdk:8-jdk-alpine} (build) & \texttt{db}    & 8080:8080 \\
  \texttt{frontend-build}   & \texttt{node:6-alpine} (build only)   & —              & — \\
  \rowcolor{rowshade}
  \texttt{frontend}         & \texttt{nginx:alpine}                 & \texttt{backend} & 80:80 \\
\end{longtable}

\vspace{0.5em}
Sample \texttt{docker-compose.yml}:

\begin{lstlisting}[language=bash, caption=docker-compose.yml]
services:
  db:
    image: mysql:5.7
    environment:
      MYSQL_ROOT_PASSWORD: <secret>
      MYSQL_DATABASE: landfill_eforms
      MYSQL_USER: appuser
      MYSQL_PASSWORD: <secret>
    volumes:
      - ./schema.sql:/docker-entrypoint-initdb.d/schema.sql

  backend:
    build: ./backend
    depends_on: [db]
    ports: ["8080:8080"]
    environment:
      SPRING_DATASOURCE_URL: jdbc:mysql://db:3306/landfill_eforms

  frontend:
    build: ./frontend    # Multi-stage: node:6-alpine build -> nginx:alpine
    depends_on: [backend]
    ports: ["80:80"]
\end{lstlisting}

\end{document}
